\section{Conclusion}
\label{sec:conclusion}

This work presented a comprehensive comparative study of eight deep learning architectures for image denoising under severe Gaussian noise ($\sigma = 100$). All models were implemented within a unified, modular framework and evaluated under identical conditions on the DIV2K dataset.

\subsection{Key Findings}

\begin{enumerate}
    \item \textbf{NAFNet-V2 achieves state-of-the-art performance} within our comparison (PSNR = 26.36~dB, SSIM = 0.747), demonstrating that activation-free blocks with Stochastic Depth and a multi-stage training pipeline (Charbonnier loss + progressive fine-tuning) are highly effective for heavy noise removal.
    
    \item \textbf{Pix2Pix with residual learning ranks second} (PSNR = 26.08~dB, SSIM = 0.733), showing that adversarial training can benefit denoising when the generator operates in residual mode and the discriminator serves as a perceptual regularizer rather than driving hallucination.
    
    \item \textbf{Residual learning is consistently beneficial:} All top-4 models (NAFNet-V2, Pix2Pix, RIDNet, DnCNN) use residual learning, confirming that predicting noise is easier than reconstructing clean images from scratch.
    
    \item \textbf{Attention mechanisms require careful design:} Bottleneck-only attention with GroupNorm outperforms all-level attention with BatchNorm by +2.36~dB, suggesting that over-application of attention gates can disrupt feature propagation in skip connections.
    
    \item \textbf{Loss function choice impacts training stability:} Charbonnier loss provides +0.28~dB over standard L1 for NAFNet, while SSIM-based fine-tuning at aggressive learning rates can destabilize pre-trained weights.
    
    \item \textbf{Training hyperparameters can dominate architecture choice:} A 10$\times$ learning rate reduction degrades DnCNN by 2.74~dB---more than the gap between the best and worst architectures. GAN training requires careful discriminator-generator LR asymmetry.
    
    \item \textbf{Skip connections dramatically accelerate convergence:} UNet variants converge in 10--36 epochs vs.\ 200 epochs for the skip-free Denoising Autoencoder, highlighting the importance of direct gradient paths.
\end{enumerate}

\subsection{Architectural Insights}

The results suggest a hierarchy of factors by importance:
\begin{enumerate}
    \item \textit{Training strategy} (multi-stage, loss scheduling) --- largest impact when combined with appropriate architecture;
    \item \textit{Learning paradigm} (residual vs.\ direct) --- critical for all architectures;
    \item \textit{Loss function} (Charbonnier, adversarial regularization) --- enables finer optimization;
    \item \textit{Multi-scale feature capture} (dilated convolutions, hierarchical encoders) --- captures noise at different frequencies;
    \item \textit{Attention mechanisms} (channel attention, simplified channel attention) --- provides adaptive feature selection;
    \item \textit{Skip connections} --- critical for training efficiency and high-frequency detail preservation.
\end{enumerate}

\subsection{Future Directions}

Potential extensions include: (1) evaluating on real-world noise distributions beyond synthetic Gaussian noise; (2) exploring ensemble methods combining the complementary strengths of RIDNet and DnCNN observed in per-image analysis; (3) investigating self-supervised denoising approaches that do not require clean reference images; and (4) extending the framework to other restoration tasks (super-resolution, deblurring, JPEG artifact removal).
